\documentclass[11pt,a4paper]{article}

\usepackage[utf8]{inputenc}
\usepackage[T1]{fontenc}
\usepackage[french]{babel}
\usepackage[margin=2.5cm]{geometry}
\usepackage{xcolor}
\usepackage{listings}
\usepackage{hyperref}
\usepackage{graphicx}

% Configuration des couleurs et du style pour le code
\definecolor{codegreen}{rgb}{0,0.6,0}
\definecolor{codegray}{rgb}{0.5,0.5,0.5}
\definecolor{codepurple}{rgb}{0.58,0,0.82}
\definecolor{backcolour}{rgb}{0.95,0.95,0.92}

\lstdefinestyle{mystyle}{
    backgroundcolor=\color{backcolour},   
    commentstyle=\color{codegreen},
    keywordstyle=\color{magenta},
    numberstyle=\tiny\color{codegray},
    stringstyle=\color{codepurple},
    basicstyle=\ttfamily\footnotesize,
    breakatwhitespace=false,         
    breaklines=true,                 
    captionpos=b,                    
    keepspaces=true,                 
    numbers=left,                    
    numbersep=5pt,                  
    showspaces=false,                
    showstringspaces=false,
    showtabs=false,                  
    tabsize=4
}
\lstset{style=mystyle}

% Definition pour GLSL
\lstdefinelanguage{GLSL}{
  morekeywords={attribute, const, uniform, varying, break, continue, do, for, while, if, else, in, out, inout, float, int, void, bool, true, false, vec2, vec3, vec4, ivec2, ivec3, ivec4, bvec2, bvec3, bvec4, mat2, mat3, mat4, sampler1D, sampler2D, sampler3D, samplerCube, struct, return, layout, location},
  morecomment=[l]{//},
  morecomment=[s]{/*}{*/},
  morestring=[b]",
}

\title{\textbf{Rapport TP — Préparation à la soutenance}}
\author{}
\date{}

\begin{document}

\maketitle

\hrule
\vspace{1em}

\section{Partie 1 — Matrices monde}

\subsection{Exercice 1.1 — Multiplication de deux matrices}

L'exercice demande d'implémenter une fonction multipliant deux matrices homogènes 4×4 stockées en colonnes (column-major). Chaque élément du résultat est la somme des produits ligne×colonne. On utilise un tableau temporaire pour permettre le cas \texttt{out == a} ou \texttt{out == b}.

\begin{lstlisting}[language=C++]
void MatrixMultiply(float* out, const float* a, const float* b)
{
    float temp[16];
    for(int c = 0; c < 4; c++) {
        for(int r = 0; r < 4; r++) {
            temp[c * 4 + r] = a[0*4 + r] * b[c*4 + 0]
                            + a[1*4 + r] * b[c*4 + 1]
                            + a[2*4 + r] * b[c*4 + 2]
                            + a[3*4 + r] * b[c*4 + 3];
        }
    }
    for(int i = 0; i < 16; i++)
        out[i] = temp[i];
}
\end{lstlisting}

Le stockage column-major signifie que l'indice \texttt{[colonne*4 + ligne]} donne l'élément à la ligne \texttt{ligne} et colonne \texttt{colonne}. La première colonne occupe les indices 0-3, la deuxième 4-7, etc.

\vspace{1em}\hrule\vspace{1em}

\subsection{Exercice 1.2 — World Matrix unique}

L'exercice demande de remplacer les matrices de transformation séparées par une unique World Matrix. L'ordre de concaténation est (de droite à gauche) :

\textbf{WorldMatrix = Translation × Rotation × Scale}

Dans notre cas, pour le dragon qui tourne autour de Y et est translaté :

\begin{lstlisting}[language=C++]
float worldDragon[16], tDragon[16], rDragon[16];
MatrixTranslation(tDragon, 4.0f, -4.0f, 0.0f);
MatrixRotationY(rDragon, time * 0.5f);
MatrixMultiply(worldDragon, tDragon, rDragon);
\end{lstlisting}

\texttt{worldDragon} est la World Matrix finale envoyée au shader en un seul appel :

\begin{lstlisting}[language=C++]
glUniformMatrix4fv(glGetUniformLocation(prog, "u_Model"), 1, GL_FALSE, worldDragon);
\end{lstlisting}

Pour le cube avec trois rotations combinées :

\begin{lstlisting}[language=C++]
float worldCube[16], tCube[16], rotX[16], rotY[16], rotZ[16];
float tmp[16], tmp2[16];
MatrixTranslation(tCube, -4.0f, 0.0f, 0.0f);
MatrixRotationY(rotY, time);
MatrixRotationX(rotX, time * 0.5f);
MatrixRotationZ(rotZ, time * 0.25f);
MatrixMultiply(tmp, rotX, rotY);
MatrixMultiply(tmp2, rotZ, tmp);
MatrixMultiply(worldCube, tCube, tmp2);
\end{lstlisting}

Cela réduit à une seule matrice uniform envoyée au GPU au lieu de plusieurs matrices séparées.

Les fonctions de base utilisées :

\begin{lstlisting}[language=C++]
void MatrixIdentity(float* m)
{
    for(int i = 0; i < 16; i++)
        m[i] = (i % 5 == 0) ? 1.0f : 0.0f;
}

void MatrixTranslation(float* m, float tx, float ty, float tz)
{
    MatrixIdentity(m);
    m[12] = tx;
    m[13] = ty;
    m[14] = tz;
}

void MatrixRotationY(float* m, float angle)
{
    MatrixIdentity(m);
    float c = cosf(angle);
    float s = sinf(angle);
    m[0] = c;  m[8] = s;
    m[2] = -s; m[10] = c;
}

void MatrixRotationX(float* m, float angle)
{
    MatrixIdentity(m);
    float c = cosf(angle);
    float s = sinf(angle);
    m[5] = c;  m[9] = -s;
    m[6] = s;  m[10] = c;
}

void MatrixRotationZ(float* m, float angle)
{
    MatrixIdentity(m);
    float c = cosf(angle);
    float s = sinf(angle);
    m[0] = c; m[4] = -s;
    m[1] = s; m[5] = c;
}
\end{lstlisting}

\vspace{1em}\hrule\vspace{1em}

\subsection{Normal Matrix (Eq 1.2)}

La normale doit être transformée par la transposée de l'inverse de la World Matrix :

\textbf{NormalMatrix = (WorldMatrix\textsuperscript{-1})\textsuperscript{T}}

La translation n'a pas de sens pour une normale (composante w = 0 en homogène). Le plus simple est de convertir la WorldMatrix en \texttt{mat3}, ce qui élimine la translation. Tant que la World Matrix ne contient pas de scale non-uniforme, \texttt{mat3(WorldMatrix)} est suffisant car pour une matrice orthogonale $M^{-1} = M^T$ donc $(M^{-1})^T = M$.

Dans le Vertex Shader :

\begin{lstlisting}[language=GLSL]
v_Normal = mat3(u_Model) * a_normal;
\end{lstlisting}

\vspace{1em}\hrule\vspace{1em}

\section{Partie 2 — Matrice vue}

\subsection{Exercice 2.1 — Fonction LookAt}

La View Matrix transforme du repère monde vers le repère de la caméra. C'est l'inverse de la World Matrix de la caméra. Comme la caméra n'a pas de scale, c'est une transformation orthogonale et on peut calculer l'inverse simplement :

\textbf{ViewMatrix = RotationMatrix\textsuperscript{T}\_{cam} × TranslationMatrix\textsuperscript{-1}\_{cam}}

La fonction LookAt prend trois paramètres : position de l'œil, position de la cible, vecteur up de référence.

Les 5 étapes de construction :

\begin{itemize}
    \item \textbf{Étape 1} — Vecteur \texttt{forward} : c'est \texttt{normalize(-(target - eye))}. Le signe moins vient du repère main droite où forward pointe hors de l'écran.
    \item \textbf{Étape 2} — Vecteur \texttt{right} : produit vectoriel \texttt{cross(up, forward)} normalisé. Règle de la main droite : up sur l'index, forward sur le majeur, right est indiqué par le pouce.
    \item \textbf{Étape 3} — Vecteur \texttt{up corrigé} : produit vectoriel \texttt{cross(forward, right)}. Forward sur le pouce, right sur l'index, up corrigé est indiqué par le majeur.
    \item \textbf{Étape 4} — Trois produits scalaires : projection de \texttt{-position} dans le repère local de la caméra. Chaque produit scalaire entre la position et un des trois vecteurs (right, up, forward) donne la composante de translation inverse.
    \item \textbf{Étape 5} — Assemblage : les 3 premières colonnes contiennent la transposée des 3 vecteurs (première colonne = toutes les composantes x, etc.), la 4ème colonne contient les produits scalaires.
\end{itemize}

\begin{lstlisting}[language=C++]
void LookAt(float* m, float ex, float ey, float ez,
                       float tx, float ty, float tz,
                       float ux, float uy, float uz)
{
    float fx = -(tx - ex);
    float fy = -(ty - ey);
    float fz = -(tz - ez);
    float fl = sqrtf(fx*fx + fy*fy + fz*fz);
    fx /= fl;
    fy /= fl;
    fz /= fl;

    float rx = uy * fz - uz * fy;
    float ry = uz * fx - ux * fz;
    float rz = ux * fy - uy * fx;
    float rl = sqrtf(rx*rx + ry*ry + rz*rz);
    rx /= rl;
    ry /= rl;
    rz /= rl;

    float uux = fy * rz - fz * ry;
    float uuy = fz * rx - fx * rz;
    float uuz = fx * ry - fy * rx;

    MatrixIdentity(m);
    m[0]  = rx;   m[4]  = ry;   m[8]  = rz;
    m[1]  = uux;  m[5]  = uuy;  m[9]  = uuz;
    m[2]  = fx;   m[6]  = fy;   m[10] = fz;

    m[12] = -(rx * ex + ry * ey + rz * ez);
    m[13] = -(uux * ex + uuy * ey + uuz * ez);
    m[14] = -(fx * ex + fy * ey + fz * ez);
}
\end{lstlisting}

La View Matrix est passée au Vertex Shader et appliquée à la position uniquement (après la World Matrix, avant la Projection). Les normales restent en espace monde pour les calculs d'illumination.

\begin{lstlisting}[language=GLSL]
void main()
{
    vec4 worldPos = u_Model * vec4(a_position, 1.0);
    v_Position  = worldPos.xyz;

    v_Normal    = mat3(u_Model) * a_normal;
    v_TexCoords = a_texcoords;

    gl_Position = u_Projection * u_View * worldPos;
}
\end{lstlisting}

La chaîne complète de transformation est donc : \texttt{gl\_Position = Projection × View × Model × vertex}

\vspace{1em}\hrule\vspace{1em}

\section{TP à rendre — Caméra orbitale}

\subsection{1. Caméra orbitale}

La caméra est positionnée sur une sphère de rayon R centrée sur la cible. On raisonne en coordonnées sphériques (R, phi, theta) contrôlées par la souris.

\textbf{Formules de conversion sphérique → cartésien} (du cours) :
\begin{itemize}
    \item $X = R \times \cos(Theta) \times \cos(Phi)$
    \item $Y = R \times \sin(Theta)$
    \item $Z = R \times \cos(Theta) \times \sin(Phi)$
\end{itemize}

\textbf{Contrôles GLFW :}
\begin{itemize}
    \item Axe horizontal de la souris → phi (azimut, rotation autour de Y)
    \item Axe vertical de la souris → theta (élévation)
    \item Molette → R (distance caméra-cible)
\end{itemize}

\textbf{Limites d'angles :} $phi \in (-\pi, +\pi)$ et $theta \in (-\pi/2, +\pi/2)$ pour éviter le retournement de la caméra.

Variables globales pour l'état de la caméra :

\begin{lstlisting}[language=C++]
static float g_Phi = 0.0f;
static float g_Theta = 0.0f;
static float g_Radius = 15.0f;

static double g_LastMouseX = 0.0;
static double g_LastMouseY = 0.0;
static bool g_MousePressed = false;

const float PI = 3.14159265359f;
\end{lstlisting}

Callback clic souris — on enregistre si le bouton gauche est pressé et la position initiale :

\begin{lstlisting}[language=C++]
void OnMouseButton(GLFWwindow* window, int button, int action, int mods)
{
    if (button == GLFW_MOUSE_BUTTON_LEFT) {
        g_MousePressed = (action == GLFW_PRESS);
        glfwGetCursorPos(window, &g_LastMouseX, &g_LastMouseY);
    }
}
\end{lstlisting}

Callback mouvement souris — on accumule les deltas dans phi et theta avec clamping :

\begin{lstlisting}[language=C++]
void OnCursorPos(GLFWwindow* window, double xpos, double ypos)
{
    if (!g_MousePressed)
        return;

    float dx = (float)(xpos - g_LastMouseX);
    float dy = (float)(ypos - g_LastMouseY);

    g_Phi   += dx * 0.005f;
    g_Theta += dy * 0.005f;

    if (g_Phi > PI)
        g_Phi -= 2.0f * PI;
    if (g_Phi < -PI)
        g_Phi += 2.0f * PI;

    if (g_Theta > PI / 2.0f - 0.01f)
        g_Theta = PI / 2.0f - 0.01f;
    if (g_Theta < -PI / 2.0f + 0.01f)
        g_Theta = -PI / 2.0f + 0.01f;

    g_LastMouseX = xpos;
    g_LastMouseY = ypos;
}
\end{lstlisting}

Callback molette — on modifie le rayon R :

\begin{lstlisting}[language=C++]
void OnScroll(GLFWwindow* win, double xoff, double yoff) {
    g_Radius -= (float)yoff;
    if (g_Radius < 1.0f) g_Radius = 1.0f;
    if (g_Radius > 50.0f) g_Radius = 50.0f;
}
\end{lstlisting}

Enregistrement des callbacks dans le main :

\begin{lstlisting}[language=C++]
glfwSetMouseButtonCallback(win, OnMouseButton);
glfwSetCursorPosCallback(win, OnCursorPos);
glfwSetScrollCallback(win, OnScroll);
\end{lstlisting}

Utilisation dans le rendu — conversion sphérique → cartésien puis appel à LookAt avec target à l'origine :

\begin{lstlisting}[language=C++]
float camX = g_Radius * cos(g_Theta) * cos(g_Phi);
float camY = g_Radius * sin(g_Theta);
float camZ = g_Radius * cos(g_Theta) * sin(g_Phi);

float view[16];
LookAt(view, camX, camY, camZ, 0, 0, 0, 0, 1, 0);
\end{lstlisting}

La position de la caméra est aussi passée au shader pour le calcul de la composante spéculaire (direction de l'observateur V) :

\begin{lstlisting}[language=C++]
glUniform3f(glGetUniformLocation(prog, "u_CameraPos"), camX, camY, camZ);
\end{lstlisting}

\vspace{1em}\hrule\vspace{1em}

\section{TP 05 — Partie 3 : Correction Gamma}

\subsection{Contexte — Le problème de la colorimétrie}

Les couleurs affichées sur un moniteur sont dites "perceptuelles". Le moniteur convertit le courant en intensité lumineuse à travers une loi de puissance (pow). Le facteur de puissance est appelé \textbf{gamma}. La plupart des moniteurs ont un gamma compris entre 1.8 et 2.4. La norme sRGB (Microsoft \& HP) fixe le gamma approximativement à \textbf{2.2}.

Cela signifie deux choses :
\begin{itemize}
    \item Les moniteurs n'ont pas une réponse linéaire : l'intensité croît en suivant une courbe de puissance
    \item Cette courbe distribue plus de valeurs aux extrémités (noirs et blancs), ce qui correspond mieux à la perception humaine
\end{itemize}

Le problème pour le programmeur : les opérations mathématiques (addition, multiplication dans les shaders) sont linéaires, mais les couleurs des pixels sont non-linéaires (compressées gamma).

\subsection{Ce qu'il faut retenir}

\begin{itemize}
    \item Le moniteur applique toujours une \textbf{décompression gamma} (mise à la puissance gamma)
    \item Une couleur visible sur un moniteur est donc toujours \textbf{non-linéaire}
    \item Les images (PNG, JPEG) stockent les couleurs avec une \textbf{compression gamma} (non-linéaire, sRGB 2.2)
    \item Une couleur compressée gamma traitée par un moniteur redevient \textbf{linéaire}
\end{itemize}

\subsection{Comment savoir si une couleur est linéaire ?}

\begin{itemize}
    \item Si la couleur a été \textbf{dessinée} (image) ou \textbf{capturée} (color picker) → elle est \textbf{non-linéaire}
    \item Si la couleur a été \textbf{générée par une équation} (normal maps, occlusion maps) → elle est \textbf{linéaire}
\end{itemize}

Les couleurs passées en uniform (matériaux, lumières) qui ont été choisies visuellement sont \textbf{non-linéaires} et doivent être linéarisées.

\vspace{1em}\hrule\vspace{1em}

\subsection{Exercice 3.1 — Correction gamma manuelle dans le Fragment Shader}

L'exercice demande d'appliquer les corrections gamma en lecture et en écriture dans le Fragment Shader.

\textbf{En lecture} — une couleur non-linéaire (texture sRGB) doit être linéarisée en appliquant la puissance 2.2. La couleur source non-linéaire = $\text{couleur}^{1/2.2}$, donc pour la linéariser on applique l'inverse $\text{couleur}^{2.2}$ :

\begin{lstlisting}[language=GLSL]
vec3 texColor = pow(texture(u_diffuseMap, v_TexCoords).rgb, vec3(2.2));
\end{lstlisting}

De même pour les couleurs de matériaux passées en uniform et choisies visuellement :

\begin{lstlisting}[language=GLSL]
vec3 matDiffuse = pow(u_material.diffuseColor, vec3(2.2));
\end{lstlisting}

\textbf{En écriture} — comme le moniteur applique automatiquement une décompression gamma, si la couleur en sortie du shader est linéaire (résultat de nos calculs d'éclairage), il faut appliquer une compression gamma de $1.0/2.2$ :

\begin{lstlisting}[language=GLSL]
vec3 finalColor = ambient + diff + spec;
FragColor = vec4(pow(finalColor, vec3(1.0/2.2)), 1.0);
\end{lstlisting}

\vspace{1em}\hrule\vspace{1em}

\subsection{Exercice 3.2 — Correction gamma automatique avec les fonctions OpenGL}

Les GPUs sont capables d'effectuer ces conversions automatiquement pour les color buffers et les textures. C'est la méthode utilisée dans notre implémentation finale.

\textbf{En lecture} — on force la conversion automatique d'une texture lors de l'échantillonnage en spécifiant un format interne sRGB. Le GPU linéarise automatiquement les texels :

\begin{lstlisting}[language=C++]
glTexImage2D(GL_TEXTURE_2D, 0, GL_SRGB8, 2, 2, 0, GL_RGB, GL_UNSIGNED_BYTE, texData);
\end{lstlisting}

Avec alpha, on utilise \texttt{GL\_SRGB8\_ALPHA8}. La composante Alpha reste toujours linéaire (séparée de SRGB8).

\textbf{En sortie} — on force les écritures dans le color buffer à convertir automatiquement de linéaire vers gamma :

\begin{lstlisting}[language=C++]
glEnable(GL_FRAMEBUFFER_SRGB);
\end{lstlisting}

Avec ces deux fonctions OpenGL, aucun \texttt{pow()} dans le shader n'est nécessaire pour les textures ni pour le framebuffer.

\textbf{Note importante} : les couleurs passées en uniform (matériaux, lumières) ne sont \textbf{pas} corrigées automatiquement par le GPU. Si elles ont été choisies visuellement, il faudrait les linéariser manuellement dans le shader ou les passer déjà linéarisées depuis le C++.

Notre implémentation dans \texttt{Initialise()} :

\begin{lstlisting}[language=C++]
glEnable(GL_FRAMEBUFFER_SRGB);

glGenTextures(1, &g_Texture);
glBindTexture(GL_TEXTURE_2D, g_Texture);
GLubyte texData[] = { 255,255,255, 0,0,0, 0,0,0, 255,255,255 };
glPixelStorei(GL_UNPACK_ALIGNMENT, 1);
glTexImage2D(GL_TEXTURE_2D, 0, GL_SRGB8, 2, 2, 0, GL_RGB, GL_UNSIGNED_BYTE, texData);
glTexParameteri(GL_TEXTURE_2D, GL_TEXTURE_MIN_FILTER, GL_NEAREST);
glTexParameteri(GL_TEXTURE_2D, GL_TEXTURE_MAG_FILTER, GL_NEAREST);
\end{lstlisting}

\vspace{1em}\hrule\vspace{1em}

\section{TP 05 — Partie 2 : Illumination ambiante hémisphérique}

\subsection{Contexte}

Dans le TP illumination précédent on a implémenté le modèle de Phong avec une composante ambiante uniforme. Ici on remplace cette ambiance uniforme par une illumination hémisphérique qui combine deux couleurs :
\begin{itemize}
    \item \textbf{SkyColor} pour l'hémisphère "up" (le ciel)
    \item \textbf{GroundColor} pour l'hémisphère "down" (le sol)
\end{itemize}

\subsection{Principe}

La technique utilise un vecteur de référence SkyDirection (typiquement \texttt{(0, 1, 0)}) et le produit scalaire entre la normale du fragment et ce vecteur.

Le produit scalaire \texttt{dot(N, SkyDirection)} est dans \texttt{[-1, +1]} quand les vecteurs sont normalisés. Pour faciliter le mixage des couleurs, on reparamètre le résultat de sorte que le domaine soit \texttt{[0, 1]} :

\textbf{HemisphereFactor = NdotSky × 0.5 + 0.5}

On utilise ensuite la fonction \texttt{mix()} du GLSL qui interpole linéairement deux couleurs :

\textbf{AmbientColor = Ia × mix(GroundColor, SkyColor, HemisphereFactor)}

\subsection{Implémentation dans le Fragment Shader}

\begin{lstlisting}[language=GLSL]
float HemisphereFactor = dot(N, normalize(u_SkyDirection)) * 0.5 + 0.5;
vec3 ambient = u_material.ambientColor * mix(u_GroundColor, u_SkyColor, HemisphereFactor) * texColor;
\end{lstlisting}

La couleur finale combine les trois composantes du modèle de Phong :

\begin{lstlisting}[language=GLSL]
vec3 FinalColor = Ambient + Diffuse + Specular;
FragColor = vec4(FinalColor, 1.0);
\end{lstlisting}

\subsection{Uniforms passés depuis le C++}

\begin{lstlisting}[language=C++]
glUniform3f(glGetUniformLocation(prog, "u_SkyDirection"), 0, 1, 0);
glUniform3f(glGetUniformLocation(prog, "u_SkyColor"), 0.6f, 0.8f, 1.0f);
glUniform3f(glGetUniformLocation(prog, "u_GroundColor"), 0.2f, 0.2f, 0.1f);
\end{lstlisting}

Le résultat est une ambiance qui varie naturellement : les faces orientées vers le haut reçoivent la couleur du ciel, celles vers le bas la couleur du sol, et celles horizontales reçoivent un mélange des deux.

\vspace{1em}\hrule\vspace{1em}

\section{Compilation et exécution}

\begin{lstlisting}[language=bash]
cd tp_final
mkdir -p build && cd build
cmake ..
make
cd .. && ./build/tp_final
\end{lstlisting}

\textbf{Contrôles :}
\begin{itemize}
    \item Clic gauche + glisser → rotation de la caméra
    \item Molette → zoom avant/arrière
\end{itemize}

\end{document}
